\section{Knowledge Graph Embeddings and Graph Neural Networks}\label{sec:fund_kge}

Section~\ref{sec:fund_kg_completion} introduced knowledge graph embedding
methods only in general terms. This section briefly introduces the two
families of representation-learning methods compared in this thesis:
translational and bilinear embedding models, which assign each entity and
relation a fixed vector learned during training, and relational graph
neural networks, which instead derive an entity's representation from its
neighbors in the graph.

The translational family is exemplified by TransE~\cite{bordes2013transe},
which models a relation as a translation from head to tail,
\[
  \mathbf{h} + \mathbf{r} \approx \mathbf{t},
\]
and is simple and efficient, but struggles with relations that are
symmetric or connect one entity to many others~\cite{wang2017kgesurvey}.
ComplEx~\cite{trouillon2016complex} addresses this by scoring a triple
with complex-valued vectors instead of a single translation, which lets
it represent both symmetric and asymmetric relations. Both methods,
however, learn a dedicated vector for every entity seen during training
and provide no representation for an entity absent from the training
graph, which is why Section~\ref{sec:fund_kg_transind} classifies them as
transductive. To score a new incident, these models therefore require an
adaptation, described in Chapter~\ref{ch:method}.
% TODO decide how TransE/ComplEx score an unseen incident: retrain,
% nearest-neighbor initialization, or exclude them from novel-incident
% evaluation entirely.

Relational graph neural networks instead compute an entity's
representation from its neighbors rather than from a fixed lookup table.
The relational graph convolutional network
(R-GCN)~\cite{schlichtkrull2018rgcn} extends the graph convolutional
network~\cite{kipf2017gcn} to multi-relational graphs, combining each
entity's representation with information from its neighbors, processed
separately per relation type, over several layers. Because this
computation draws on an entity's actual neighbors rather than a
pre-trained vector, it can in principle be applied to an entity that was
never seen during training -- provided its initial representation does
not depend on its identity (e.g.\ a vector of known attribute values
rather than a learned per-entity embedding).
% TODO decide R-GCN initialization of new incident nodes: attribute
% vector, zero vector, or learned type embedding shared across incidents.

These three mechanisms -- the rule-based exact-match baseline of
Section~\ref{sec:fund_bmw_lookup}, translational and bilinear embeddings,
and relational graph neural networks -- form the three model families
compared in this thesis (Chapter~\ref{ch:method}). They trade off
generalization and interpretability: the rule-based baseline is fully
transparent but generalizes only to an identical historical match;
embedding models generalize further through geometric similarity, at the
cost of a score that is harder to translate into checkable evidence; and
graph neural networks generalize furthest, to incidents never observed
before, but their layered message passing makes it least obvious which
part of the neighborhood drove a prediction. This trade-off is exactly
what the main research question (Section~\ref{sec:intro_rq}) asks about,
and it motivates evaluating accuracy and faithfulness jointly, across
novelty levels, rather than in isolation.
